% Partie : set-theory
% Chapitre : spine
% Section : stagesbasics

\documentclass[../../../include/open-logic-section]{subfiles}

\begin{document}

\olfileid[fr]{sth}{spine}{Valphabasic}
\olsection{Propriétés élémentaires des étapes}

Pour faire ressortir l'importance fondamentale de la définition des
$V_\alpha$, nous allons présenter quelques résultats élémentaires à
leur sujet. Commençons par une définition :\footnote{Il n'existe pas
de terminologie consacrée pour «~potent~» ; c'est le terme employé
par \citet{ButtonLT1}.}
\begin{defn}
L'ensemble $A$ est \emph{potent} si et seulement si
$\forall x((\exists y \in A)x \subseteq y \lif x \in A)$.
\end{defn}
\begin{lem}\ollabel{Valphabasicprops}
Pour chaque ordinal $\alpha$ :
\begin{enumerate}
	\item\ollabel{Valphatrans} $V_\alpha$ est transitif.
	\item\ollabel{Valphapotent} $V_\alpha$ est potent.
	\item\ollabel{Valphacum} Si $\gamma \in \alpha$, alors
	$V_\gamma \in V_\alpha$ (et donc aussi
	$V_\gamma \subseteq V_\alpha$ par \olref{Valphatrans}).
\end{enumerate}
\end{lem}

\begin{proof}
Nous démontrons ces affirmations simultanément par récurrence
transfinie. Supposons, comme hypothèse de récurrence, que
\olref{Valphatrans}--\olref{Valphacum} valent pour tout ordinal
$\beta < \alpha$.

Le cas $\alpha = \emptyset$ est immédiat.

Supposons $\alpha = \ordsucc{\beta}$. Pour établir
\olref{Valphacum}, si $\gamma \in \alpha$, alors
$V_\gamma \subseteq V_\beta$,
par l'hypothèse de récurrence
ou par égalité,
et donc
$V_\gamma \in \Pow{V_\beta} = V_\alpha$. Pour établir
\olref{Valphapotent}, supposons $A \subseteq B \in V_\alpha$,
c'est-à-dire $A \subseteq B \subseteq V_\beta$. Alors
$A \subseteq V_\beta$, donc $A \in V_\alpha$. Pour établir
\olref{Valphatrans}, remarquons que, si $x \in A \in V_\alpha$,
alors $A \subseteq V_\beta$, donc $x \in V_\beta$. Comme
$V_\beta$ est transitif par hypothèse, $x \subseteq V_\beta$ et,
par conséquent, $x \in V_\alpha$.

Supposons que $\alpha$ soit un ordinal limite. Pour établir
\olref{Valphacum}, si $\gamma \in \alpha$, alors
$\gamma \in \ordsucc{\gamma} \in \alpha$. Par hypothèse,
$V_\gamma \in V_{\ordsucc{\gamma}}$ et donc
$V_\gamma \in \bigcup_{\beta \in \alpha} V_\beta = V_\alpha$.
Pour établir \olref{Valphatrans} et \olref{Valphapotent}, il suffit
d'observer qu'une union d'ensembles transitifs, ou respectivement
potents, est transitive, ou respectivement potente.
\end{proof}

\begin{lem}\ollabel{Valphanotref}
Pour chaque ordinal $\alpha$, $V_\alpha \notin V_\alpha$.
\end{lem}

\begin{proof}
Procédons par récurrence transfinie. Il est évident que
$V_\emptyset \notin V_\emptyset$.

Si $V_{\ordsucc{\alpha}} \in V_{\ordsucc{\alpha}} =
\Pow{V_\alpha}$, alors
$V_{\ordsucc{\alpha}} \subseteq V_\alpha$. Or
$V_\alpha \in V_{\ordsucc{\alpha}}$ par le
\olref{Valphabasicprops}, d'où $V_\alpha \in V_\alpha$.
Par contraposée, si $V_\alpha \notin V_\alpha$, alors
$V_{\ordsucc{\alpha}} \notin V_{\ordsucc{\alpha}}$.

Si $\alpha$ est un ordinal limite et si
$V_\alpha \in V_\alpha = \bigcup_{\beta \in \alpha}V_\beta$,
alors $V_\alpha \in V_\beta$ pour un certain $\beta \in \alpha$.
Mais nous avons aussi $V_\beta \in V_\alpha$ ; par le
\olref{Valphabasicprops} appliqué deux fois, il s'ensuit que
$V_\beta \in V_\beta$. Par contraposée, si
$V_\beta \notin V_\beta$ pour tout $\beta \in \alpha$, alors
$V_\alpha \notin V_\alpha$.
\end{proof}

\begin{cor}
Pour tous ordinaux $\alpha, \beta$ :
$\alpha \in \beta$ si et seulement si $V_\alpha \in V_\beta$.
\end{cor}

\begin{proof}
Le \olref{Valphabasicprops} donne un sens de l'équivalence.
Réciproquement, supposons $V_\alpha \in V_\beta$. Alors
$\alpha \neq \beta$ par le \olref{Valphanotref}. De plus,
$\beta \notin \alpha$ : sinon, nous aurions
$V_\beta \in V_\alpha$, puis $V_\beta \in V_\beta$ par le
\olref{Valphabasicprops} appliqué deux fois, en contradiction avec
le \olref{Valphanotref}. Donc $\alpha \in \beta$ par trichotomie.
\end{proof}
\noindent
Tout cela permet de voir en $V_\alpha$ l'étape d'indice $\alpha$ de
la hiérarchie. En voici la raison.

Nous pouvons assurément concevoir la formation des $V_\alpha$
comme un processus \emph{itératif} : les ordinaux rendent compte de
l'itération. De plus, si une étape se forme avant une autre,
c'est-à-dire si $V_\beta \in V_\alpha$, donc si
$\beta \in \alpha$, le processus est \emph{cumulatif}, puisque
$V_\beta \subseteq V_\alpha$. Enfin, nous formons bien \emph{toutes}
les collections possibles d'ensembles disponibles à une étape
antérieure : toute étape successeur $V_{\ordsucc{\alpha}}$ est
l'ensemble des parties de sa devancière $V_\alpha$.

Bref, dans $\ZFminus$, nous avons \emph{presque} entièrement précisé
notre conception de la hiérarchie cumulative et itérative des
ensembles. (Il reste évidemment à justifier le remplacement.)

\end{document}
